% Options for packages loaded elsewhere
\PassOptionsToPackage{unicode}{hyperref}
\PassOptionsToPackage{hyphens}{url}
\PassOptionsToPackage{dvipsnames,svgnames,x11names}{xcolor}
%
\documentclass[
  letterpaper,
  DIV=11,
  numbers=noendperiod]{scrartcl}

\usepackage{amsmath,amssymb}
\usepackage{iftex}
\ifPDFTeX
  \usepackage[T1]{fontenc}
  \usepackage[utf8]{inputenc}
  \usepackage{textcomp} % provide euro and other symbols
\else % if luatex or xetex
  \usepackage{unicode-math}
  \defaultfontfeatures{Scale=MatchLowercase}
  \defaultfontfeatures[\rmfamily]{Ligatures=TeX,Scale=1}
\fi
\usepackage{lmodern}
\ifPDFTeX\else  
    % xetex/luatex font selection
\fi
% Use upquote if available, for straight quotes in verbatim environments
\IfFileExists{upquote.sty}{\usepackage{upquote}}{}
\IfFileExists{microtype.sty}{% use microtype if available
  \usepackage[]{microtype}
  \UseMicrotypeSet[protrusion]{basicmath} % disable protrusion for tt fonts
}{}
\makeatletter
\@ifundefined{KOMAClassName}{% if non-KOMA class
  \IfFileExists{parskip.sty}{%
    \usepackage{parskip}
  }{% else
    \setlength{\parindent}{0pt}
    \setlength{\parskip}{6pt plus 2pt minus 1pt}}
}{% if KOMA class
  \KOMAoptions{parskip=half}}
\makeatother
\usepackage{xcolor}
\ifLuaTeX
  \usepackage{luacolor}
  \usepackage[soul]{lua-ul}
\else
  \usepackage{soul}
  
\fi
\setlength{\emergencystretch}{3em} % prevent overfull lines
\setcounter{secnumdepth}{-\maxdimen} % remove section numbering
% Make \paragraph and \subparagraph free-standing
\makeatletter
\ifx\paragraph\undefined\else
  \let\oldparagraph\paragraph
  \renewcommand{\paragraph}{
    \@ifstar
      \xxxParagraphStar
      \xxxParagraphNoStar
  }
  \newcommand{\xxxParagraphStar}[1]{\oldparagraph*{#1}\mbox{}}
  \newcommand{\xxxParagraphNoStar}[1]{\oldparagraph{#1}\mbox{}}
\fi
\ifx\subparagraph\undefined\else
  \let\oldsubparagraph\subparagraph
  \renewcommand{\subparagraph}{
    \@ifstar
      \xxxSubParagraphStar
      \xxxSubParagraphNoStar
  }
  \newcommand{\xxxSubParagraphStar}[1]{\oldsubparagraph*{#1}\mbox{}}
  \newcommand{\xxxSubParagraphNoStar}[1]{\oldsubparagraph{#1}\mbox{}}
\fi
\makeatother


\providecommand{\tightlist}{%
  \setlength{\itemsep}{0pt}\setlength{\parskip}{0pt}}\usepackage{longtable,booktabs,array}
\usepackage{calc} % for calculating minipage widths
% Correct order of tables after \paragraph or \subparagraph
\usepackage{etoolbox}
\makeatletter
\patchcmd\longtable{\par}{\if@noskipsec\mbox{}\fi\par}{}{}
\makeatother
% Allow footnotes in longtable head/foot
\IfFileExists{footnotehyper.sty}{\usepackage{footnotehyper}}{\usepackage{footnote}}
\makesavenoteenv{longtable}
\usepackage{graphicx}
\makeatletter
\def\maxwidth{\ifdim\Gin@nat@width>\linewidth\linewidth\else\Gin@nat@width\fi}
\def\maxheight{\ifdim\Gin@nat@height>\textheight\textheight\else\Gin@nat@height\fi}
\makeatother
% Scale images if necessary, so that they will not overflow the page
% margins by default, and it is still possible to overwrite the defaults
% using explicit options in \includegraphics[width, height, ...]{}
\setkeys{Gin}{width=\maxwidth,height=\maxheight,keepaspectratio}
% Set default figure placement to htbp
\makeatletter
\def\fps@figure{htbp}
\makeatother

\KOMAoption{captions}{tableheading}
\makeatletter
\@ifpackageloaded{tcolorbox}{}{\usepackage[skins,breakable]{tcolorbox}}
\@ifpackageloaded{fontawesome5}{}{\usepackage{fontawesome5}}
\definecolor{quarto-callout-color}{HTML}{909090}
\definecolor{quarto-callout-note-color}{HTML}{0758E5}
\definecolor{quarto-callout-important-color}{HTML}{CC1914}
\definecolor{quarto-callout-warning-color}{HTML}{EB9113}
\definecolor{quarto-callout-tip-color}{HTML}{00A047}
\definecolor{quarto-callout-caution-color}{HTML}{FC5300}
\definecolor{quarto-callout-color-frame}{HTML}{acacac}
\definecolor{quarto-callout-note-color-frame}{HTML}{4582ec}
\definecolor{quarto-callout-important-color-frame}{HTML}{d9534f}
\definecolor{quarto-callout-warning-color-frame}{HTML}{f0ad4e}
\definecolor{quarto-callout-tip-color-frame}{HTML}{02b875}
\definecolor{quarto-callout-caution-color-frame}{HTML}{fd7e14}
\makeatother
\makeatletter
\@ifpackageloaded{caption}{}{\usepackage{caption}}
\AtBeginDocument{%
\ifdefined\contentsname
  \renewcommand*\contentsname{Table of contents}
\else
  \newcommand\contentsname{Table of contents}
\fi
\ifdefined\listfigurename
  \renewcommand*\listfigurename{List of Figures}
\else
  \newcommand\listfigurename{List of Figures}
\fi
\ifdefined\listtablename
  \renewcommand*\listtablename{List of Tables}
\else
  \newcommand\listtablename{List of Tables}
\fi
\ifdefined\figurename
  \renewcommand*\figurename{Figure}
\else
  \newcommand\figurename{Figure}
\fi
\ifdefined\tablename
  \renewcommand*\tablename{Table}
\else
  \newcommand\tablename{Table}
\fi
}
\@ifpackageloaded{float}{}{\usepackage{float}}
\floatstyle{ruled}
\@ifundefined{c@chapter}{\newfloat{codelisting}{h}{lop}}{\newfloat{codelisting}{h}{lop}[chapter]}
\floatname{codelisting}{Listing}
\newcommand*\listoflistings{\listof{codelisting}{List of Listings}}
\makeatother
\makeatletter
\makeatother
\makeatletter
\@ifpackageloaded{caption}{}{\usepackage{caption}}
\@ifpackageloaded{subcaption}{}{\usepackage{subcaption}}
\makeatother

\ifLuaTeX
  \usepackage{selnolig}  % disable illegal ligatures
\fi
\usepackage{bookmark}

\IfFileExists{xurl.sty}{\usepackage{xurl}}{} % add URL line breaks if available
\urlstyle{same} % disable monospaced font for URLs
\hypersetup{
  pdftitle={BSTA 551 Syllabus},
  colorlinks=true,
  linkcolor={blue},
  filecolor={Maroon},
  citecolor={Blue},
  urlcolor={Blue},
  pdfcreator={LaTeX via pandoc}}


\title{BSTA 551 Syllabus}
\author{}
\date{}

\begin{document}
\maketitle


\subsection{Key Course Info}\label{key-course-info}

\begin{itemize}
\tightlist
\item
  \textbf{If an assignment on Sakai is closed} or you are submitting
  late work, please email me AND the TA your work
\item
  The in-person class instruction will end on Wednesday, March 11, 2026.
  \textbf{All coursework is expected to be completed by Thursday of
  finals week, March 19, 2026 at 11pm.}
\end{itemize}

\subsection{Description}\label{description}

Welcome to BSTA 551! In this course we will study the theoretical
foundaton of statistical inference which includes estimation and
hypothesis testing. This is the official course description:

This course introduces theoretical foundation in Biostatistics. Topics
will include distributions of random variables (location-scale families
and exponential families), data reduction (sufficiency and completeness
of statistics), methods of estimation (method of moment estimators and
maximum likelihood estimators), convergence, finite and large sample
properties of estimators, interval estimation, hypothesis testing,
asymptotic tests (likelihood-ratio tests, score tests, and Wald tests),
and statistical simulations to evaluate statistical methods.

\subsubsection{Course Learning
Objectives}\label{course-learning-objectives}

At the end of this course, students should be able to\ldots{}

\begin{enumerate}
\def\labelenumi{\arabic{enumi}.}
\tightlist
\item
  Explain the major concepts and theorems in statistical inference.
\item
  Connect theoretical concepts to statistical analyses.
\item
  Conduct simulations to study and evaluate statistical methods.
\end{enumerate}

\subsection{Instructors}\label{instructors}

\href{instructors.qmd}{Here is the instructor page.}

\subsection{Meeting Times}\label{meeting-times}

Mondays~~~~~~~~~~10:00 AM -- 12:00 PM PST (see Sakai for room)

Wednesdays~~~~10:00 AM -- 12:00 PM PST

\subsubsection{Known Exceptions}\label{known-exceptions}

\begin{itemize}
\tightlist
\item
  Monday, January 19: No class (holiday)
\item
  Wednesday, February 11: No class
\item
  Monday, February 16: No class (holiday)
\item
  March 16 and 18: Finals week, no in person teaching, possibly virtual
  office hours as needed
\end{itemize}

\subsection{Materials}\label{materials}

\subsubsection{Textbook}\label{textbook}

\paragraph{Modern Mathematical Statistics with Applications, 3rd
ed.}\label{modern-mathematical-statistics-with-applications-3rd-ed.}

\begin{itemize}
\item
  \textbf{Author:} JL Devore, KN Berk, MA Carlton
\item
  \href{https://librarysearch.ohsu.edu/permalink/01ALLIANCE_OHSU/okd4j/alma99900609916301858}{Textbook
  available online through library}
\item
  Citation: Devore JL, Berk KN, Carlton MA. Modern Mathematical
  Statistics with Applications. Third edition. Springer; 2021.
  doi:10.1007/978-3-030-55156-8
\item
  Focus on chapters 6-10
\end{itemize}

\paragraph{Supplemental Readings
(Optional)}\label{supplemental-readings-optional}

\begin{itemize}
\item
  \emph{Statistical Inference}, Casella and Berger, 2nd ed.~(This was
  the previous textbook BSTA 551-552 Math Stat., very theoretical)
  \href{https://librarysearch.ohsu.edu/permalink/01ALLIANCE_OHSU/19jn9i0/alma99900537344401858}{ebook
  OHSU library link}
\item
  \emph{An Introduction to R}
  (\href{http://cran.r-project.org/manuals.html}{free pdf available})
\item
  \emph{Statistical Inference via Data Science, A ModernDive into R and
  the Tidyverse}, Ismay, Kim, Valdivia, 2nd
  ed.~\href{https://moderndive.com/v2/index.html}{free ebook online}
\end{itemize}

\subsubsection{Online Resources}\label{online-resources}

\paragraph{Sakai}\label{sakai}

While most course materials will be delivered online through this
website, assignments will be turned in through
\href{https://sakai.ohsu.edu/}{Sakai}, OHSU's course management system.
I will include a link on this website to the Sakai assignment page.~

\paragraph{PennState STAT 415 Website}\label{pennstate-stat-415-website}

Dr.~Wakim linked to Penn State's probability course materials for 550.
They also have the subsequent inference course on the website as well.
They have all their \href{https://online.stat.psu.edu/stat415/}{course
notes posted on this page} if you wish to have an alternative reference.

\paragraph{University of South Carolina STAT 713
Website}\label{university-of-south-carolina-stat-713-website}

This is a more thoeretical version (similar to previous math/stats
courses taught at OHSU) of statistical inference taught by Professor
Tebbs at U of SC. Their notes
\href{https://people.stat.sc.edu/Tebbs/stat713/index.htm}{which are
linked on the course website} are very well written. If you want more
theoretical context for what we are learning in this course this is a
good source. Note, it uses Casella \& Berger's textbook
\emph{Statistical Inference} as a reference.

\paragraph{R: Statistical Computing
Software}\label{r-statistical-computing-software}

R/Rstudio will be used to complete some homework assignments. Please see
\href{https://nwakim.github.io/BSTA_550_25F/syllabus.html}{Dr.~Wakim's
BSTA 550 syllabus} section on R if you need any references or refreshers
on how to use.

\subsection{Assessment}\label{assessment}

\subsubsection{Breakdown}\label{breakdown}

50\% Homework (Weekly)

20\% Midterm

20\% Final

10\% Attendence (Two exit tickets a week)

\paragraph{Grading \& Requirements}\label{grading-requirements}

Letter grades will be assigned roughly according to the following
scheme: A (\textgreater=93\%), A- (90-92\%), B+ (88-89\%), B(83-87\%),
B- (82-80\%), C+(78-79\%), C(73-77\%), C- (70-72\%), D (60 -- 69\%),
F(\textless60\%).

\paragraph{Homework grading}\label{homework-grading}

{[}Identical to Dr.~Nicky Wakim's BSTA 550 approach to homework!{]}

No student has the same amount of time available to dedicate to
homework. This class may not be a priority to you, you may be taking
several other courses, or you may need to dedicate time to other
activities. Homework assignments are \textbf{formative assessments},
meaning its purpose is to help you learn and practice. To reduce the
pressure on you to have perfect homework (the first time around), I have
a very simple grading policy: \textbf{Your homework will be given a
check mark if you turn in 75\% of the question parts \emph{completed}
(whether the 75\% is correct or wrong).} I highly encourage you to stay
up-to-date with the homeworks and put in as much effort as you can.
\textbf{This will be the most helpful work in this class!}

If you turn in the homework on time, I will give you feedback (on one or
more complete problems). There is no penalty for turning in the homework
late, \textbf{but you will not get feedback on your work.} Please make
sure to check the solutions or go to office hours to assess your work.

\paragraph{Viewing Grades in Sakai}\label{viewing-grades-in-sakai}

Points you receive for graded activities will be posted to the Sakai
Gradebook. Click on the Gradebook link on the left navigation to view
your points.

\subsection{Course \& Instructor
Evaluations}\label{course-instructor-evaluations}

\subsubsection{Ongoing Course Feedback}\label{ongoing-course-feedback}

\href{https://forms.office.com/r/2yhdKvbkNX}{Submit feedback throughout
the course here.}

\subsubsection{Final Course Feedback}\label{final-course-feedback}

At the conclusion of the course, you will be asked to complete a formal
online review of the course and the instructor. Your feedback on this
University evaluation is critical to improving future student learning
in this course as well as providing metrics relevant to the instructor's
career advancement (or lack of). Since our class is on the smaller side,
everyone's participation is needed for feedback to be released.

\subsection{Schedule}\label{schedule}

\href{./schedule.qmd}{Please refer to the Schedule page.} I will make
changes to this schedule if we need more or less time on a concept. You
do not need to read the corresponding chapters in the textbook for each
class.

\subsubsection{Assignments}\label{assignments}

Written problem sets, approximately weekly, usually assigned Wednesday
and due the following Thursday.

Homework problems will be posted on the course website. Solutions should
be posted as a pdf in the corresponding Sakai submission box.

\subsection{Course Policies and
Resources}\label{course-policies-and-resources}

\subsubsection{Late Work Policy}\label{late-work-policy}

I encourage you to make your best effort to submit all assignments on
time, but I understand circumstances arise that are beyond our control.
Please see this
\href{https://myuni.swansea.ac.uk/academic-life/extenuating-circumstances/\#how-can-the-university-support-me-and-what-should-i-do=is-expanded&what-are-extenuating-circumstances=is-expanded&what-isnt-an-extenuating-circumstance=is-expanded&what-supporting-evidence-is-required=is-expanded&who-should-i-contact-to-make-an-extenuating-circumstances-application=is-expanded&who-will-make-a-decision-on-my-application=is-expanded}{Swansea
University's page on extenuating circumstances} for some examples. Not
all circumstances are covered here, so please reach out if you have
questions.~

\begin{itemize}
\item
  The class will end on Wednesday, March 18, 2026. \textbf{All
  coursework is expected to be completed by March 19, 2026 at 11pm.} If
  you have \emph{extenuating circumstances}, and need additional time to
  complete class assignments, please contact me. Together, we will come
  up with a plan for completion and to sort out registrar logistics. If
  you need an Incomplete, generally this requires request two weeks
  before the end of class.
\item
  If you have extenuating circumstances that may jeopardize your ability
  to do work for several weeks, please contact me. We will come up with
  a plan to keep you on track in the course and prevent any delay in
  your education.
\item
  For homework, you will have TWO no-questions-asked, 3-day extensions:
  one for the first assignment part and one for either the solutions or
  presentation. Please use this wisely! You just need to send me a quick
  email saying ``I am using my no-questions-asked extension for Homework
  \_\_ assignment/solutions/presentation.''
\item
  For homework, I ask you to email me directly about any late
  submissions. You can explain your circumstances and may ask us for an
  extension. I am very likely to grant an extension, but \textbf{I want
  to emphasize how important it will be to stay on track with your
  homework}! Your group is depending on you, and delaying homework may
  only add stress on the next homework!
\item
  If you have a emergency involving your self, family, pet, friend,
  classmate, or anything/one deemed important to you, \textbf{please do
  not worry about immediately contacting me}. We can work something out
  after your emergency. If I contact you during an emergency, it is only
  because I am worried, and you do NOT need to respond until you are
  able.~
\end{itemize}

\subsubsection{Regrade Policy}\label{regrade-policy}

If you think a question was incorrectly graded, first compare your
answer to the answer key. If you believe a re-grade would be
appropriate, write an email to me containing the question and a short
explanation as to why the question(s) was/were incorrectly graded.
Deadline: One week after assignments were returned to class (late
requests will not be considered).

\subsubsection{Attendance Policy}\label{attendance-policy}

You are expected to attend class, participate in-class polls, and
complete the exit ticket. For students who miss class or need a review,
I will make video and audio recordings of lectures available. There are
no guarantees against technical or other challenges for the recording
availability or quality.

\textbf{You will need to attend all classes.} There are 17 classes
total, so you are welcome to watch the recordings or come in-person.
While I want attendance to be a flexible thing, I need to set certain
requirements around in-person attendance to align with the school's
policy. Attendance is measured through exit tickets that will be due 7
days after each class.

This is meant to keep you on track within the course and prevent a pile
up of material. Make sure to complete the exit ticket at the end of
class to demonstrate attendance.

\subsubsection{Plagiarism and
Attribution}\label{plagiarism-and-attribution}

Please note that this section is directly sourced from Dr.~Nicky Wakim's
Course Policies for BSTA 550 which has itself has been motivated by
\href{https://dmice.ohsu.edu/bedricks/courses/bmi525/policies.html}{Dr.~Steven
Bedrick's Course Policies and Grading site} for BMI 525. (Note that this
is a good example of informal attribution of someone else's work.)

In this class, it is easy to use ChatGPT or other AI tools to solve your
homework for you. Many problems follow a basic structure that is
especially easy for ChatGPT to solve. In this class, you may use ChatGPT
to help with your homework. You may even ask for direct answers.
However, there are a few things I do not want you to do:

\begin{itemize}
\item
  \textbf{Do not copy ChatGPT's answer directly into your homework.}
  Your homework is graded for full credit if you turn it in, in any
  state, so turning in ChatGPT's answers is unacceptable. I rather see
  half-written answers that show what you're thinking than see a correct
  answer from ChatGPT.
\item
  \textbf{Do not stop once ChatGPT answered a question.} If it gives an
  explanation, interact with it! Make sure you understand the thought
  process of ChatGPT. Try writing out the process to help cement it in
  your head. Check the answer with what we learn in class.
\item
  \textbf{Do not use ChatGPT on our exams!} Hence, you need to really
  understand how to solve these problems even if you use ChatGPT on the
  homework.
\end{itemize}

At the end of the day, ChatGPT is a resource that will be available to
you in a job and outside of school. Thus, we should use it as a tool in
school as well! Let me know if ChatGPT helped you understand something!
I would love to incorporate it into future classes!

\begin{tcolorbox}[enhanced jigsaw, colbacktitle=quarto-callout-important-color!10!white, bottomrule=.15mm, title=\textcolor{quarto-callout-important-color}{\faExclamation}\hspace{0.5em}{Important}, breakable, colframe=quarto-callout-important-color-frame, toprule=.15mm, bottomtitle=1mm, colback=white, leftrule=.75mm, coltitle=black, arc=.35mm, titlerule=0mm, opacityback=0, opacitybacktitle=0.6, rightrule=.15mm, left=2mm, toptitle=1mm]

You can think of this class as assembling a toolbox. When a handyperson
starts working for the first time, they need to buy their tools. For
their first few jobs, they might need help finding their tools or
remembering which tool is best used for what action. Eventually, they
get to know their tools well, and using them appropriately becomes
second nature.

\textbf{For now, ChatGPT can help us find and use our tools, but we need
to work towards using them as second nature!}

\end{tcolorbox}

\subsection{Course Expectations}\label{course-expectations}

\subsubsection{Instructor Expectations}\label{instructor-expectations}

\ul{\emph{Commitment to your learning and your success\\
}}I believe that everyone has the ability to be successful in this
course and I have put a lot of effort into designing the course in a way
that maximizes your learning to ensure your success. Please talk to me
before or after class or stop by my office if there is anything you want
to discuss or about which you are unclear. I want to be supportive of
your learning and growth.

\ul{\emph{Inclusive \& supportive learning community}}\\
I believe that learning happens best when we all learn together, as a
community. This means creating a space characterized by generous
listening, civility, humility, patience, and hospitality. I will attempt
to promote a safe climate where we examine content from multiple
perspectives. I will strive to create and maintain a classroom
atmosphere in which you feel free to both listen to others and express
your views and ask questions to increase your learning.

\ul{\emph{Openness to feedback}}\\
I appreciate straightforward feedback from you regarding how well the
class is meeting your needs. Let me know if material is not clear or
when its relevance to the student learning outcomes for the course is
not apparent. In particular, let me know if you identify bias or
stereotyping in my teaching materials as I will seek to continuously
improve. Please also let me know if there's an aspect of the class you
find particularly interesting, helpful, or enjoyable!

\ul{\emph{Responsiveness\\
}}I will monitor email as well as the discussion board daily and try
respond to all messages within 24 hours Monday-Friday.

\ul{\emph{Clear guidelines and prompt feedback on assignments}}\\
I will provide clear instructions for all assignments, and a grading
rubric when applicable. The TAs and I will provide detailed feedback on
your submissions and will update grades promptly in Sakai.

\subsubsection{Student Expectations and
Resources}\label{student-expectations-and-resources}

\ul{\emph{Attend class}}\textbf{\hfill\break
}You are expected to attend at least 10 scheduled class meetings
in-person. Attendance is taken through exit tickets. If you have issues
accessing the poll on a specific day, please let me know.~

\ul{\emph{Participate\textbf{\hfill\break
}}}I encourage you to participate actively in class. I will expect all
students, and all instructors, to be respectful of each other's
contributions, whether I agree with them or not. Professional
interactions are expected.

\ul{\emph{Build rapport}}\\
If you find that you have any trouble keeping up with assignments or
other aspects of the course, make sure you let me know as early as
possible. As you will find, building rapport and effective relationships
are key to becoming an effective professional. Make sure that you are
proactive in informing me when difficulties arise during the quarter so
that I can help you find a solution in regards to coursework.

\ul{\emph{Complete assignments}}\\
All assignments for this course will be submitted electronically through
Sakai unless otherwise instructed. I encourage you to make your best
effort to submit all assignments on time, but I understand that
sometimes circumstances arise that are beyond our control. If you need
an extension, please contact me in congruence with the Late Policy.

\ul{\emph{Seek help if you need it}}\textbf{\hfill\break
}I believe it is important to support the physical and emotional
well‐being of my students. If you are experiencing physical and/or
mental health issues, I encourage you to use the resources on campus
such as those listed below. If you have a health issue that is affecting
your performance or participation in the course, and/or if you need help
connecting with these resources, please contact me.

\begin{itemize}
\item
  Student Health and Wellness Center (SHW),
  \href{https://www.ohsu.edu/education/student-health-and-wellness}{Website},
  503-494-8665 (OHSU Students only)
\item
  Student Health and Counseling (SHAC),
  \href{https://www.pdx.edu/health-counseling/}{Website}, 503-725-2800
\end{itemize}

\ul{\emph{Inform your instructor of any accommodations needed}}\\
You should speak with or email me before or during the first week of
classes regarding any special needs. Students seeking academic
accommodations should register with the appropriate service under the
School policies below.

Some religious holidays may occur on regularly scheduled class days.
Because available class hours are so limited in number, we will have to
hold class on all such days. Class video recordings will be available
and you are encouraged to engage with the material outside of the
regular class time. Please email me about your absence. I will excuse
the absece from your grade. You are also encouraged to come to office
hours with questions from the session.

\ul{\emph{Commit to integrity}}\textbf{\hfill\break
}As a student in this course (and at PSU or OHSU) you are expected to
maintain high degrees of professionalism, commitment to active learning
and participation in this class and also integrity in your behavior in
and out of the classroom.

Cheating and other forms of academic misconduct will not be tolerated in
this course and will be dealt with firmly. Student academic misconduct
refers to behavior that includes plagiarism, cheating on assignments,
fabrication of data, falsification of records or official documents,
intentional misuse of equipment or materials (including library
materials), or aiding and abetting the perpetration of such acts.
Preparation of exams, assigned on an individual basis, must represent
each student's own individual effort. When used, resource materials
should be cited in conventional reference format.

\subsection{Course Communications}\label{course-communications}

\textbf{Sakai announcements\\
}For important/urgent matters, I will communicate with you using
announcements via Sakai that will be delivered to your OHSU Email
account as well as displayed in the Sakai course site Announcements
section.

\textbf{E-mail}\\
E-mail should be used only for messages that are private in nature.
Please send private messages to my OHSU email address
(minnier@ohsu.edu).

\subsection{Further Student Resources}\label{further-student-resources}

\subsubsection{Academic Success Center}\label{academic-success-center}

OHSU houses an Academic Success Center for all students. Their mission
is to create a center for learning support where ALL learners can
discover the resources and community that they need for finding academic
success at OHSU. They provide many services to students, including:
learning skills support, writing support, English for speakers of other
languages (ESOL) support, and individual and group content tutoring.
\href{https://ohsuitg.sharepoint.com/sites/AcademicSuccessCenter}{Check
out the SharePoint site for the Academic Success Center.}

\subsubsection{Student Wellness}\label{student-wellness}

I am committed to supporting the physical and emotional well-being of my
students. Both PSU and OHSU have designated centers for student health.
For OHSU, students can visit the
\href{https://www.ohsu.edu/education/behavioral-health}{Behavioral
Health site}, where you can find more information including the number
to make an appointment. \textbf{All student visits are free.} OHSU
students also have access to PSU's
\href{https://www.pdx.edu/health-counseling/counseling}{Counseling
Services} through the school's Student Health \& Counseling.~Information
on additional student resources for OHSU students are available on the
OHSU \href{https://ohsu-psu-sph.org/health_and_wellness/}{Health and
Wellness Resource page}.~

\subsubsection{Support for Food
Insecurity}\label{support-for-food-insecurity}

Students across the country experience food insecurity at high rates.
OHSU and PSU both provide a list of resources to help combat food
insecurity. Of note, the Committee to Improve Student Food Security
(CISFS) at PSU provides a
\href{https://www.pdx.edu/student-access-center/free-food-market}{Free
Food Market} on the second Monday of each month. OHSU also provides
\href{https://o2.ohsu.edu/student-central/health-wellness/student-food-resources/snap.cfm}{SNAP
Enrollment Assistance}. The
\href{https://oregonhunger.org/snap-for-students/}{Supplemental
Nutrition Assistance Program (SNAP)} allocates money towards food for
individuals below a certain income level. If you make less than \$2,430
monthly, you may wish to enroll.

\subsubsection{Support for Students with
Children}\label{support-for-students-with-children}

Students who have children can use the PSU resource:
\href{https://www.pdx.edu/students-with-children/}{Resource Center for
Students with Children}. Resources are mostly focused on students with
younger children. There are several great resources available,
including: family-friendly study spaces, new baby starter packs, free
kids clothing, and further information on financial resources for
childcare.

\section{Additional information}\label{additional-information}

Please go to the Syllabus tab in Sakai for the following information:

\begin{itemize}
\item
  School Policies and Resources
\item
  OHSU Competencies
\item
  Institutional Policies and Resources
\end{itemize}

I do not include them here because I do not have the bandwidth to make
sure all links and resources are up-to-date.




\end{document}
